\section{Quy trình nghiệp vụ}

Mỗi quy trình dưới đây được trình bày theo các bước tuần tự, có ghi rõ tác nhân thực hiện (actor), dữ liệu vào/ra, điều kiện chuyển bước và các trường hợp ngoại lệ. Bộ quy trình này là cơ sở trực tiếp để xây dựng các sơ đồ BPMN/flowchart và đặc tả use-case ở các chương thiết kế tiếp theo của đồ án.

\subsection{Đăng ký và xác thực nhà bán, kiểm duyệt chứng nhận}

\begin{enumerate}
    \item Nhà bán (nông dân hoặc hợp tác xã) đăng ký tài khoản, nhập thông tin định danh (căn cước công dân hoặc mã số thuế đối với hợp tác xã) và thông tin trang trại, vùng canh tác.
    \item Hệ thống gửi mã xác thực (OTP) qua số điện thoại hoặc email; nếu nhập sai quá 5 lần, tài khoản bị khóa tạm thời trong 30 phút.
    \item Nhà bán tải lên giấy chứng nhận VietGAP, GlobalGAP hoặc hữu cơ (nếu có), kèm mã số chứng nhận, cơ quan cấp và ngày hết hạn.
    \item Bộ phận kiểm duyệt đối chiếu thông tin chứng nhận với dữ liệu của cơ quan cấp (nếu có kết nối) hoặc xác minh thủ công.
    \item Hệ thống ghi nhận kết quả: nếu đạt, tài khoản chuyển trạng thái “Đã xác thực” và được gắn nhãn chứng nhận tương ứng hiển thị công khai; nếu không đạt, hệ thống gửi thông báo lý do và hướng dẫn bổ sung.
\end{enumerate}

Ngoại lệ: khi chứng nhận còn hiệu lực dưới 30 ngày, hệ thống tự động nhắc nhà bán gia hạn; nếu hết hạn mà chưa gia hạn, nhãn chứng nhận trên sản phẩm liên quan sẽ tự động bị ẩn cho đến khi được cập nhật, sản phẩm không bị gỡ bỏ. Trường hợp phát hiện chứng nhận giả mạo qua xác minh hoặc khiếu nại, tài khoản bị khóa vĩnh viễn, được ghi vào danh sách vi phạm và không được phép đăng ký lại bằng cùng thông tin định danh.

\begin{figure}[H]
\centering
\begin{adjustbox}{max totalsize={\textwidth}{\dimexpr\textheight-4.6cm\relax},center}
\begin{tikzpicture}[farm/chart]
    \node[farm/terminal] (start) at (0,0) {Bắt đầu};
    \node[farm/process] (register) at (0,-1.25) {Nhà bán: đăng ký, khai báo định danh và thông tin trang trại};
    \node[farm/process] (otp) at (0,-2.75) {Hệ thống gửi OTP;\\nhà bán nhập mã xác thực};
    \node[farm/decision] (valid) at (0,-4.5) {OTP hợp lệ?};
    \node[farm/decision] (attempts) at (6,-4.5) {Sai quá 5 lần?};
    \node[farm/exception] (lock) at (6,-6.4) {Hệ thống: khóa tạm thời\\30 phút};
    \node[farm/process] (ekyc) at (0,-6.75) {Nhà bán: xác minh danh tính điện tử (eKYC) --- ảnh giấy tờ định danh và ảnh chân dung};
    \node[farm/process] (evidence) at (0,-8.55) {Nhà bán: tải chứng nhận (nếu có), mã số, cơ quan cấp và hạn dùng};
    \node[farm/process] (review) at (0,-10.4) {Kiểm duyệt: đối chiếu kết quả eKYC và chứng nhận với cơ quan cấp hoặc kiểm tra thủ công};
    \node[farm/decision] (approved) at (0,-12.45) {eKYC và hồ sơ\\đạt yêu cầu?};
    \node[farm/processnarrow] (supplement) at (6,-12.45) {Hệ thống: thông báo lý do của bước không đạt; nhà bán bổ sung};
    \node[farm/process] (verified) at (0,-14.75) {Hệ thống: trạng thái ``Đã xác thực'', kích hoạt quyền đăng bán; gắn nhãn chứng nhận nếu có};
    \node[farm/process] (monitor) at (0,-16.45) {Hệ thống: theo dõi hiệu lực chứng nhận trong quá trình hoạt động};
    \node[farm/processnarrow] (expiry) at (6,-16.45) {Còn dưới 30 ngày: nhắc gia hạn.\\Hết hạn: ẩn nhãn đến khi cập nhật; không gỡ sản phẩm.};
    \node[farm/terminal] (finish) at (0,-18.1) {Hoàn tất đăng ký};
    \node[farm/note] at (2.7,-19.9) {\textbf{Ngoại lệ xuyên suốt:} nếu phát hiện chứng nhận giả mạo qua xác minh hoặc khiếu nại, khóa vĩnh viễn tài khoản, ghi vào danh sách vi phạm và chặn đăng ký lại bằng cùng thông tin định danh.};
    \foreach \source/\target in {start/register,register/otp,otp/valid,ekyc/evidence,evidence/review,review/approved,verified/monitor,monitor/finish}
        \draw[farm/arrow] (\source) -- (\target);
    \draw[farm/arrow] (valid) -- node[farm/label,right] {Có} (ekyc);
    \draw[farm/arrow] (valid) -- node[farm/label,above] {Không} (attempts);
    \draw[farm/arrow] (attempts.north) |- node[farm/label,near start,right] {Không: nhập lại} (otp.east);
    \draw[farm/arrow] (attempts) -- node[farm/label,right] {Có} (lock);
    \draw[line width=0.7pt, draw=black!75] (lock.east) -- (8.6,-6.4) -- node[farm/label,right] {Sau 30 phút} (8.6,-2.75) -- (6,-2.75);
    \draw[farm/arrow] (approved) -- node[farm/label,right] {Có} (verified);
    \draw[farm/arrow] (approved) -- node[farm/label,above] {Không} (supplement);
    \draw[farm/arrow] (supplement.west) -- (3.3,-12.45) -- (3.3,-6.75) -- (ekyc.east);
    \draw[farm/arrow,dashed] (monitor) -- (expiry);
\end{tikzpicture}
\end{adjustbox}
\caption{Lưu đồ quy trình đăng ký, xác thực nhà bán và kiểm duyệt chứng nhận}
\label{fig:flow-dang-ky}
\end{figure}

\subsection{Đăng và kiểm duyệt sản phẩm, sinh mã QR truy xuất}

\begin{enumerate}
    \item Nhà bán đã xác thực tạo sản phẩm mới với đầy đủ thông tin: tên, mô tả, hình ảnh, đơn vị tính, giá bán lẻ, giá bán sỉ theo bậc (nếu có), sản lượng và ngày dự kiến thu hoạch.
    \item Nhà bán ghi nhật ký canh tác theo từng giai đoạn (gieo trồng, chăm sóc, thu hoạch), mỗi lần ghi gồm ngày tháng, hoạt động và hình ảnh minh chứng nếu có; các bản ghi được gắn với một lô sản phẩm (batch) cụ thể.
    \item Hệ thống kiểm duyệt tự động sơ bộ: kiểm tra hình ảnh không trùng lặp, mô tả không chứa từ khóa cấm, giá không bất thường so với mặt bằng chung của danh mục.
    \item Bộ phận kiểm duyệt thực hiện kiểm tra thủ công lần cuối đối với sản phẩm mới hoặc nhà bán chưa có lịch sử giao dịch; nhà bán có uy tín lâu năm có thể áp dụng cơ chế hậu kiểm để rút ngắn thời gian chờ duyệt.
    \item Hệ thống sinh mã QR truy xuất duy nhất cho từng lô sản phẩm, liên kết tới trang truy xuất công khai gồm hồ sơ nhà bán, chứng nhận, nhật ký canh tác và các mốc vận chuyển liên quan.
\end{enumerate}

Điều kiện: sản phẩm chỉ được duyệt và hiển thị công khai khi có tối thiểu một bản ghi nhật ký canh tác gần nhất và hình ảnh sản phẩm thực tế. Ngoại lệ: khi nhà bán thay đổi thông tin trọng yếu sau khi đã duyệt (giá thay đổi vượt 20\% hoặc mô tả chất lượng thay đổi), sản phẩm được chuyển về trạng thái chờ duyệt lại; các thay đổi nhỏ (còn hàng/hết hàng) không yêu cầu duyệt lại.

\clearpage
\begin{figure}[p]
\centering
\begin{tikzpicture}[farm/chart]
    \node[farm/terminal] (start) at (0,0) {Nhà bán đã xác thực};
    \node[farm/process] (create) at (0,-1.5) {Nhà bán: tạo sản phẩm, khai báo giá lẻ/sỉ, sản lượng, ngày thu hoạch và nhật ký theo lô};
    \node[farm/decision] (complete) at (0,-3.6) {Có nhật ký gần nhất\\và ảnh thực tế?};
    \node[farm/side] (supplement) at (6,-3.6) {Nhà bán: bổ sung nhật ký và hình ảnh còn thiếu};
    \node[farm/process] (automatic) at (0,-5.6) {Hệ thống: kiểm tra ảnh trùng lặp, từ khóa cấm và giá bất thường};
    \node[farm/decision] (passed) at (0,-7.4) {Đạt kiểm tra\\sơ bộ?};
    \node[farm/side] (revise) at (6,-7.4) {Chưa duyệt; nhà bán chỉnh sửa thông tin sản phẩm};
    \node[farm/decision] (trusted) at (0,-9.5) {Được áp dụng\\hậu kiểm?};
    \node[farm/side] (manual) at (6,-9.5) {Kiểm duyệt: kiểm tra thủ công sản phẩm mới / nhà bán chưa có lịch sử giao dịch};
    \node[farm/decision] (accepted) at (6,-11.8) {Được duyệt?};
    \node[farm/process] (publish) at (0,-11.8) {Hệ thống: duyệt, công khai sản phẩm; sinh QR riêng cho lô, liên kết hồ sơ, chứng nhận, nhật ký và vận chuyển};
    \node[farm/decision] (changed) at (0,-14) {Thay đổi\\trọng yếu?};
    \node[farm/side] (pending) at (6,-14) {Giá thay đổi vượt 20\% hoặc đổi mô tả chất lượng:\\chuyển chờ duyệt lại};
    \node[farm/terminal] (finish) at (0,-16) {Tiếp tục hiển thị};
    \node[farm/connector] (reviseout) at (9,-7.4) {A};
    \node[farm/connector] (revisein) at (6,-0.2) {A};
    \node[farm/connector] (reviewout) at (9,-14) {B};
    \node[farm/connector] (reviewin) at (6,-5.6) {B};
    \node[farm/note] at (2.7,-17.6) {\textbf{Cập nhật nhỏ:} thay đổi còn hàng/hết hàng không cần duyệt lại. Nhánh hậu kiểm chỉ áp dụng cho nhà bán uy tín theo cơ chế được cho phép; việc kiểm tra được thực hiện sau khi công khai sản phẩm.\\[2pt]\textbf{Điểm nối:} các vòng tròn cùng ký hiệu A hoặc B nối tiếp cùng một luồng.};
    \foreach \source/\target in {start/create,create/complete,automatic/passed,manual/accepted,publish/changed}
        \draw[farm/arrow] (\source) -- (\target);
    \draw[farm/arrow] (complete) -- node[farm/label,right] {Có} (automatic);
    \draw[farm/arrow] (complete) -- node[farm/label,above] {Không} (supplement);
    \draw[farm/arrow] (supplement.north) |- (create.east);
    \draw[farm/arrow] (passed) -- node[farm/label,right] {Có} (trusted);
    \draw[farm/arrow] (passed) -- node[farm/label,above] {Không} (revise);
    \draw[farm/arrow] (revise) -- (reviseout);
    \draw[farm/arrow] (revisein.south) |- (create.east);
    \draw[farm/arrow] (trusted) -- node[farm/label,right] {Có} (publish);
    \draw[farm/arrow] (trusted) -- node[farm/label,above] {Không} (manual);
    \draw[farm/arrow] (accepted) -- node[farm/label,above] {Có} (publish);
    \draw[farm/arrow] (accepted.east) -- (8.6,-11.8) -- (8.6,-8.5) -| node[farm/label,near start,right] {Không} (revise.south);
    \draw[farm/arrow] (changed) -- node[farm/label,right] {Không} (finish);
    \draw[farm/arrow] (changed) -- node[farm/label,above] {Có} (pending);
    \draw[farm/arrow] (pending) -- (reviewout);
    \draw[farm/arrow] (reviewin) -- (automatic);
\end{tikzpicture}
\caption{Flowchart đăng và kiểm duyệt sản phẩm, sinh mã QR truy xuất.}
\label{fig:flow-san-pham}
\end{figure}
\clearpage

\subsection{Quy trình mua lẻ (B2C)}

\begin{enumerate}
    \item Người tiêu dùng tìm kiếm và lọc sản phẩm theo danh mục, vùng trồng, chứng nhận hoặc mức giá.
    \item Người tiêu dùng xem chi tiết sản phẩm, có thể tra cứu trang truy xuất nguồn gốc trước khi quyết định mua.
    \item Người tiêu dùng thêm sản phẩm vào giỏ hàng với số lượng giới hạn theo sản lượng còn lại của lô hàng.
    \item Người tiêu dùng thanh toán qua ví điện tử, chuyển khoản hoặc thanh toán khi nhận hàng (COD), đồng thời chọn khung giờ giao phù hợp với hàng tươi.
    \item Hệ thống xác nhận đơn hàng, tạm trừ sản lượng khỏi lô hàng và thông báo cho nhà bán chuẩn bị hàng.
    \item Nhà bán đóng gói theo tiêu chuẩn dành cho hàng dễ hư hỏng và bàn giao cho đơn vị vận chuyển.
    \item Người tiêu dùng xác nhận đã nhận hàng và có thể đánh giá sản phẩm, nhà bán trong khoảng thời gian quy định sau khi nhận hàng.
\end{enumerate}

Ngoại lệ: nếu thanh toán trực tuyến thất bại, đơn hàng được giữ ở trạng thái chờ thanh toán tối đa 15 phút trước khi tự hủy và giải phóng lại sản lượng. Nếu người nhận từ chối nhận hàng COD, đơn chuyển trạng thái hoàn hàng và chi phí vận chuyển hai chiều được xử lý theo chính sách đã công bố. Khi nhiều đơn đặt cùng lúc vượt sản lượng còn lại, hệ thống ưu tiên đơn thanh toán/xác nhận trước; các đơn còn lại được tự động hủy kèm thông báo hết hàng.

\clearpage
\begin{figure}[p]
\centering
\begin{tikzpicture}[farm/chart]
    \node[farm/terminal] (start) at (0,0) {Bắt đầu mua lẻ};
    \node[farm/process] (browse) at (0,-1.35) {Người mua: tìm kiếm, lọc, xem chi tiết và truy xuất nguồn gốc};
    \node[farm/process] (cart) at (0,-2.85) {Người mua: thêm vào giỏ theo sản lượng còn lại; chọn khung giờ giao và phương thức thanh toán};
    \node[farm/decision] (payment) at (0,-5) {Thanh toán thành công\\hoặc chọn COD?};
    \node[farm/side] (wait) at (6,-5) {Thanh toán trực tuyến thất bại:\\giữ đơn chờ tối đa 15 phút, cho phép thanh toán lại};
    \node[farm/decision] (retry) at (6,-7.2) {Thành công\\trong 15 phút?};
    \node[farm/decision] (stock) at (0,-7.2) {Còn đủ\\sản lượng?};
    \node[farm/exception] (cancel) at (6,-9.4) {Hệ thống: hủy đơn, thông báo lý do; giải phóng sản lượng đã giữ nếu có};
    \node[farm/process] (confirm) at (0,-9.4) {Hệ thống: xác nhận đơn, tạm trừ sản lượng và báo nhà bán};
    \node[farm/process] (ship) at (0,-11) {Nhà bán: đóng gói hàng dễ hư hỏng, bàn giao vận chuyển};
    \node[farm/decision] (refused) at (0,-12.9) {Người nhận\\từ chối COD?};
    \node[farm/exception] (return) at (6,-12.9) {Hoàn hàng; xử lý phí vận chuyển hai chiều theo chính sách};
    \node[farm/process] (receive) at (0,-14.9) {Người mua: xác nhận nhận hàng, đánh giá sản phẩm và nhà bán trong thời hạn quy định};
    \node[farm/terminal] (finish) at (0,-16.5) {Kết thúc};
    \node[farm/note] at (2.7,-18.2) {\textbf{Đặt đồng thời:} ưu tiên đơn thanh toán/xác nhận trước khi phân bổ sản lượng; đơn còn lại bị hủy và thông báo hết hàng.\\[2pt]\textbf{Hàng hư hỏng:} xử lý theo flowchart vận chuyển và khiếu nại tại Hình~\ref{fig:flow-van-chuyen}.};
    \foreach \source/\target in {start/browse,browse/cart,cart/payment,wait/retry,confirm/ship,ship/refused,receive/finish}
        \draw[farm/arrow] (\source) -- (\target);
    \draw[farm/arrow] (payment) -- node[farm/label,right] {Có} (stock);
    \draw[farm/arrow] (payment) -- node[farm/label,above] {Không} (wait);
    \draw[farm/arrow] (retry) -- node[farm/label,above] {Có} (stock);
    \draw[farm/arrow] (retry) -- node[farm/label,right] {Không} (cancel);
    \draw[farm/arrow] (stock) -- node[farm/label,right] {Có} (confirm);
    \draw[farm/arrow] (stock.south east) |- (3.2,-8.5) |- node[farm/label,near end,above] {Không} (cancel.west);
    \draw[farm/arrow] (refused) -- node[farm/label,right] {Không} (receive);
    \draw[farm/arrow] (refused) -- node[farm/label,above] {Có} (return);
    \draw[farm/arrow] (return.south) |- (finish.east);
    \draw[farm/arrow] (cancel.east) -- (8.5,-9.4) -- (8.5,-16.5) -- (finish.east);
\end{tikzpicture}
\caption{Flowchart mua lẻ (B2C), thanh toán và xử lý đơn hàng.}
\label{fig:flow-mua-le}
\end{figure}
\clearpage

\subsection{Quy trình mua sỉ (B2B)}

\begin{enumerate}
    \item Doanh nghiệp thu mua gửi yêu cầu báo giá (RFQ) tới một hoặc nhiều nhà bán, nêu rõ sản phẩm, số lượng dự kiến, tiêu chuẩn chất lượng yêu cầu và lịch giao mong muốn.
    \item Hệ thống kiểm tra số lượng yêu cầu có đạt số lượng tối thiểu (MOQ) của nhà bán hay không; nếu chưa đạt, hệ thống gợi ý gộp đơn với doanh nghiệp khác cùng khu vực hoặc chuyển hướng sang kênh mua lẻ.
    \item Nhà bán phản hồi báo giá theo giá bậc (ví dụ: mức giá khác nhau theo từng ngưỡng số lượng), thời gian có thể đáp ứng và điều kiện thanh toán.
    \item Hai bên thương lượng qua hệ thống nhắn tin nội bộ cho đến khi thống nhất số lượng, đơn giá và lịch giao.
    \item Hệ thống sinh hợp đồng điện tử từ mẫu chuẩn, điền đầy đủ các điều khoản đã thống nhất, bao gồm điều khoản phạt vi phạm và điều khoản thanh toán/công nợ.
    \item Hai bên ký xác nhận hợp đồng bằng chữ ký điện tử hoặc mã xác nhận trong hệ thống.
    \item Nhà bán giao hàng theo từng đợt theo lịch đã ký; mỗi đợt giao, hệ thống sinh hóa đơn tương ứng và cập nhật vào sổ công nợ của doanh nghiệp thu mua.
    \item Doanh nghiệp thu mua xác nhận nhận hàng và chất lượng từng đợt, sau đó thanh toán theo kỳ hạn công nợ đã thỏa thuận.
\end{enumerate}

Ngoại lệ: khi nhà bán không giao đủ số lượng hoặc đúng lịch một đợt giao, điều khoản phạt trong hợp đồng được áp dụng và sự việc được ghi nhận vào lịch sử uy tín của nhà bán. Khi doanh nghiệp thu mua chậm thanh toán quá hạn công nợ, hệ thống nhắc nhở tự động và có thể tạm khóa quyền đặt yêu cầu báo giá mới cho đến khi hoàn tất thanh toán. Trường hợp một trong hai bên muốn hủy hợp đồng giữa kỳ, việc xử lý tuân theo điều khoản hủy đã ký kết từ đầu.

\clearpage
\begin{figure}[p]
\centering
\begin{tikzpicture}[farm/chart]
    \node[farm/terminal] (start) at (0,0) {Bắt đầu mua sỉ};
    \node[farm/process] (rfq) at (0,-1.4) {Doanh nghiệp: gửi RFQ gồm sản phẩm, số lượng, chất lượng và lịch giao mong muốn};
    \node[farm/decision] (moq) at (0,-3.3) {Đạt MOQ\\của nhà bán?};
    \node[farm/side] (alternative) at (6,-3.3) {Hệ thống: gợi ý gộp đơn cùng khu vực hoặc chuyển mua lẻ};
    \node[farm/terminal] (retail) at (6,-5.3) {Chuyển quy trình B2C};
    \node[farm/process] (quote) at (0,-5.3) {Nhà bán: báo giá theo bậc, thời gian đáp ứng và điều kiện thanh toán};
    \node[farm/process] (negotiate) at (0,-6.8) {Hai bên: thương lượng qua hệ thống nhắn tin nội bộ};
    \node[farm/decision] (agreed) at (0,-8.6) {Thống nhất\\điều khoản?};
    \node[farm/process] (contract) at (0,-10.6) {Hệ thống: sinh hợp đồng, gồm phạt vi phạm và công nợ; hai bên ký điện tử / mã xác nhận};
    \node[farm/process] (delivery) at (0,-12.4) {Nhà bán giao theo từng đợt; hệ thống lập hóa đơn và cập nhật sổ công nợ};
    \node[farm/process] (settle) at (0,-14.1) {Doanh nghiệp: xác nhận số lượng, chất lượng từng đợt; thanh toán theo kỳ hạn};
    \node[farm/decision] (complete) at (0,-16) {Còn đợt giao\\theo hợp đồng?};
    \node[farm/side] (reconcile) at (6,-16) {Hai bên: đối soát, hoàn tất thanh toán công nợ theo kỳ hạn};
    \node[farm/terminal] (finish) at (0,-18) {Hoàn tất hợp đồng};
    \node[farm/exception,text width=3.8cm] (exceptions) at (6,-10.6) {\textbf{Ngoại lệ hợp đồng}\\[3pt]Giao thiếu / trễ: áp dụng phạt và ghi lịch sử uy tín.\\[4pt]Quá hạn công nợ: nhắc tự động; có thể khóa RFQ mới đến khi trả đủ.\\[4pt]Hủy giữa kỳ: xử lý theo điều khoản hủy đã ký.};
    \foreach \source/\target in {start/rfq,rfq/moq,quote/negotiate,negotiate/agreed,contract/delivery,delivery/settle,settle/complete}
        \draw[farm/arrow] (\source) -- (\target);
    \draw[farm/arrow] (moq) -- node[farm/label,right] {Có} (quote);
    \draw[farm/arrow] (moq) -- node[farm/label,above] {Không} (alternative);
    \draw[farm/arrow] (alternative.north) |- node[farm/label,near start,right] {Gộp đơn, gửi lại} (rfq.east);
    \draw[farm/arrow] (alternative) -- node[farm/label,right] {Chọn mua lẻ} (retail);
    \draw[farm/arrow] (agreed) -- node[farm/label,right] {Có} (contract);
    \draw[farm/arrow] (agreed.west) -- (-3.5,-8.6) |- node[farm/label,near start,left] {Chưa} (quote.west);
    \draw[farm/arrow] (complete) -- node[farm/label,above] {Không} (reconcile);
    \draw[farm/arrow] (reconcile.south) |- (finish.east);
    \draw[farm/arrow] (complete.west) -- (-3.5,-16) |- node[farm/label,near start,left] {Có} (delivery.west);
    \draw[farm/arrow,dashed] (contract) -- (exceptions);
\end{tikzpicture}
\caption{Flowchart mua sỉ (B2B): báo giá, hợp đồng, giao hàng và công nợ.}
\label{fig:flow-mua-si}
\end{figure}
\clearpage

\subsection{Vận chuyển hàng dễ hư hỏng và xử lý khiếu nại}

\begin{enumerate}
    \item Nhà bán đóng gói theo hướng dẫn bảo quản riêng cho từng nhóm nông sản (nhiệt độ, độ ẩm, vật liệu chống dập) do hệ thống gợi ý theo danh mục sản phẩm.
    \item Hệ thống ưu tiên lựa chọn phương thức vận chuyển có hỗ trợ bảo quản lạnh hoặc thời gian giao nhanh đối với các danh mục được đánh dấu dễ hư hỏng; đơn hàng có thời gian giao vượt ngưỡng an toàn sẽ được cảnh báo trước khi khách xác nhận đặt hàng.
    \item Đơn vị vận chuyển cập nhật trạng thái đơn hàng theo thời gian thực trong suốt quá trình giao nhận.
    \item Người nhận kiểm tra hàng ngay khi nhận; nếu phát hiện hư hỏng, gửi khiếu nại trong khung thời gian quy định kèm hình ảnh hoặc video minh chứng.
    \item Hệ thống phân loại khiếu nại theo nguyên nhân khả năng (đóng gói/chất lượng từ nhà bán, lỗi vận chuyển, hoặc bảo quản sai từ phía người nhận).
    \item Bộ phận phụ trách ra quyết định xử lý: hoàn tiền toàn phần, hoàn tiền một phần, đổi lô hàng mới, hoặc từ chối khiếu nại nếu bằng chứng không đầy đủ.
\end{enumerate}

Ngoại lệ: khiếu nại lặp lại nhiều lần từ cùng một nhà bán trong thời gian ngắn sẽ tự động được gắn cờ cảnh báo và tạm hạ mức ưu tiên hiển thị sản phẩm cho đến khi được rà soát; khiếu nại có dấu hiệu bất thường từ phía người mua sẽ được chuyển cho bộ phận quản trị xem xét thủ công trước khi hoàn tiền.

\clearpage
\begin{figure}[p]
\centering
\begin{adjustbox}{max totalsize={\textwidth}{\dimexpr\textheight-2.4cm\relax},center}
\begin{tikzpicture}[farm/chart]
    \node[farm/terminal] (start) at (0,0) {Đơn hàng dự kiến};
    \node[farm/process] (method) at (0,-1.35) {Hệ thống: ưu tiên giao nhanh / bảo quản lạnh theo danh mục dễ hư hỏng};
    \node[farm/decision] (unsafe) at (0,-3.2) {Thời gian giao vượt ngưỡng an toàn?};
    \node[farm/processnarrow] (warn) at (6,-3.2) {Hệ thống: cảnh báo trước khi khách xác nhận đặt hàng};
    \node[farm/process] (confirm) at (0,-5.6) {Khách xác nhận đặt hàng;\\nhà bán đóng gói theo nhiệt độ, độ ẩm và vật liệu phù hợp};
    \node[farm/process] (transport) at (0,-7.3) {Đơn vị vận chuyển: giao hàng, cập nhật trạng thái theo thời gian thực};
    \node[farm/decision] (damage) at (0,-9.4) {Người nhận kiểm tra: hàng hư hỏng?};
    \node[farm/terminal] (received) at (6,-9.4) {Nhận hàng, hoàn tất};
    \node[farm/process] (claim) at (0,-11.4) {Người nhận: khiếu nại trong thời hạn, kèm ảnh / video minh chứng};
    \node[farm/decision] (suspicious) at (0,-13.4) {Người mua có\\dấu hiệu bất thường?};
    \node[farm/processnarrow] (manual) at (6,-13.4) {Quản trị: xem xét thủ công trước khi quyết định hoàn tiền};
    \node[farm/process] (classify) at (0,-15.8) {Hệ thống: phân loại nguyên nhân từ nhà bán, vận chuyển hoặc bảo quản của người nhận};
    \node[farm/process] (resolve) at (0,-17.7) {Bộ phận phụ trách: hoàn tiền toàn phần / một phần, đổi lô mới hoặc từ chối nếu thiếu bằng chứng};
    \node[farm/terminal] (finish) at (0,-19.4) {Thông báo kết quả xử lý};
    \node[farm/exception] (repeated) at (6,-17.6) {Khiếu nại lặp lại từ cùng nhà bán trong thời gian ngắn:\\gắn cờ và tạm hạ ưu tiên hiển thị đến khi rà soát};
    \foreach \source/\target in {start/method,method/unsafe,confirm/transport,transport/damage,claim/suspicious,classify/resolve,resolve/finish}
        \draw[farm/arrow] (\source) -- (\target);
    \draw[farm/arrow] (unsafe) -- node[farm/label,right] {Không} (confirm);
    \draw[farm/arrow] (unsafe) -- node[farm/label,above] {Có} (warn);
    \draw[farm/arrow] (warn.south) |- (confirm.east);
    \draw[farm/arrow] (damage) -- node[farm/label,right] {Có} (claim);
    \draw[farm/arrow] (damage) -- node[farm/label,above] {Không} (received);
    \draw[farm/arrow] (suspicious) -- node[farm/label,right] {Không} (classify);
    \draw[farm/arrow] (suspicious) -- node[farm/label,above] {Có} (manual);
    \draw[farm/arrow] (manual.south) |- (classify.east);
    \draw[farm/arrow,dashed] (resolve) -- (repeated);
\end{tikzpicture}
\end{adjustbox}
\caption{Flowchart vận chuyển hàng dễ hư hỏng và xử lý khiếu nại.}
\label{fig:flow-van-chuyen}
\end{figure}
\clearpage

\subsection{Quản trị hệ thống}

\begin{enumerate}
    \item Quản trị viên thực hiện kiểm duyệt định kỳ hồ sơ nhà bán mới, sản phẩm mới hoặc có thay đổi trọng yếu, và các chứng nhận sắp hết hạn.
    \item Quản trị viên giám sát các cảnh báo tự động từ hệ thống, bao gồm tỷ lệ khiếu nại bất thường và các dấu hiệu giao dịch có nguy cơ gian lận.
    \item Vi phạm được xử lý theo quy trình phân cấp từ nhắc nhở, cảnh cáo công khai, tạm khóa cho đến khóa vĩnh viễn tài khoản, tùy theo mức độ và số lần vi phạm, có lưu vết đầy đủ lý do và quyết định.
    \item Quản trị viên truy xuất báo cáo thống kê định kỳ để hỗ trợ ra quyết định vận hành và báo cáo cho các bên liên quan.
\end{enumerate}

Ngoại lệ: khi phát hiện vi phạm nghiêm trọng như giả mạo chứng nhận hoặc gian lận, tài khoản bị khóa ngay lập tức mà không qua các bước nhắc nhở trung gian, đồng thời toàn bộ đơn hàng đang xử lý của tài khoản đó được rà soát để bảo vệ quyền lợi của khách hàng đã đặt trước.

\begin{figure}[H]
\centering
\begin{adjustbox}{max totalsize={\textwidth}{\dimexpr\textheight-4.6cm\relax},center}
\begin{tikzpicture}[farm/chart]
    \node[farm/terminal] (start) at (0,0) {Bắt đầu đợt giám sát};
    \node[farm/process] (review) at (0,-1.5) {Quản trị: kiểm duyệt hồ sơ mới, sản phẩm mới / thay đổi trọng yếu và chứng nhận sắp hết hạn};
    \node[farm/process] (monitor) at (0,-3.2) {Quản trị: theo dõi cảnh báo tỷ lệ khiếu nại bất thường và giao dịch có nguy cơ gian lận};
    \node[farm/decision] (violation) at (0,-5.2) {Phát hiện\\vi phạm?};
    \node[farm/decision] (severe) at (0,-7.9) {Vi phạm\\nghiêm trọng?};
    \node[farm/exception] (lock) at (6,-7.9) {Giả mạo chứng nhận / gian lận:\\khóa ngay tài khoản, không qua nhắc nhở trung gian};
    \node[farm/process] (sanction) at (0,-10.5) {Quản trị: xử lý theo mức độ và số lần vi phạm: nhắc nhở, cảnh cáo công khai, tạm khóa hoặc khóa vĩnh viễn};
    \node[farm/exception] (orders) at (6,-10.5) {Rà soát toàn bộ đơn đang xử lý của tài khoản, bảo vệ khách đã đặt trước};
    \node[farm/process] (audit) at (0,-12.7) {Hệ thống / quản trị: lưu vết đầy đủ lý do và quyết định xử lý};
    \node[farm/process] (report) at (0,-14.5) {Quản trị: truy xuất báo cáo định kỳ, hỗ trợ quyết định vận hành và báo cáo các bên liên quan};
    \node[farm/terminal] (finish) at (0,-16.2) {Kết thúc đợt giám sát};
    \node[farm/note] at (2.7,-18) {\textbf{Nguyên tắc phân cấp:} các hình thức xử lý được lựa chọn theo mức độ và số lần vi phạm, không bắt buộc đi qua mọi mức. Quy trình giám sát được thực hiện định kỳ; trường hợp nghiêm trọng được xử lý ngay khi phát hiện.};
    \foreach \source/\target in {start/review,review/monitor,monitor/violation,lock/orders,sanction/audit,audit/report,report/finish}
        \draw[farm/arrow] (\source) -- (\target);
    \draw[farm/arrow] (violation) -- node[farm/label,right] {Có} (severe);
    \draw[farm/arrow] (violation.west) -- (-3.6,-5.2) |- node[farm/label,near start,left] {Không} (report.west);
    \draw[farm/arrow] (severe) -- node[farm/label,right] {Không} (sanction);
    \draw[farm/arrow] (severe) -- node[farm/label,above] {Có} (lock);
    \draw[farm/arrow] (orders.south) |- (audit.east);
\end{tikzpicture}
\end{adjustbox}
\caption{Lưu đồ quy trình quản trị hệ thống, xử lý vi phạm và báo cáo}
\label{fig:flow-quan-tri}
\end{figure}